\documentclass[11pt]{article}
\usepackage[margin=1in]{geometry}
\usepackage{amsmath,amssymb,amsthm}
\usepackage{graphicx}
\usepackage[hidelinks]{hyperref}
\emergencystretch=3em

\newtheorem*{theorem*}{Theorem}

\title{A Positive Answer to Question 4.6 of arXiv:2602.12699}
\author{Run fa\_banach\_001, lane 6}
\date{June 28, 2026}

\begin{document}
\maketitle

\section*{Status}

This packet gives a candidate full positive answer to Question 4.6 of Basu,
Behera, and Seal, arXiv:2602.12699. The question asks whether semi SCS points
are stable under binary $\ell_\infty$-sums in the direction not covered by
Theorem 4.5 of the source paper.

\section*{Source Question}

The source paper defines a point $x_0$ of a bounded set $C$ to be a
\emph{semi Small Combination of Slice point} (semi SCS point) if, for every
$\varepsilon>0$, there is a convex combination of slices of $C$ contained in
$B(x_0,\varepsilon)$.

The relevant source context is Theorem 4.5 and Question 4.6 on page 8:

\begin{center}
\includegraphics[width=0.95\linewidth]{figures/open_problem_crop.png}
\end{center}

In text form, Question 4.6 asks:

\begin{quote}
Let $X$ and $Y$ be Banach spaces. If $x\in \operatorname{semiSCS}(B_X)$ and
$y\in \operatorname{semiSCS}(B_Y)$, does it follow that
$(x,y)\in \operatorname{semiSCS}(B_{X\oplus_\infty Y})$?
\end{quote}

\section*{Idea of the Proof}

Take convex combinations of slices around $x$ and around $y$. For each pair
of component slices, the average functional
\[
  (u,v)\mapsto \frac12 f_i(u)+\frac12 g_j(v)
\]
defines a slice of the $\ell_\infty$-sum ball. If its aperture is small, then
membership in this product slice forces the first coordinate into the
$i$th slice of $B_X$ and the second coordinate into the $j$th slice of $B_Y$.
Using product weights $\lambda_i\mu_j$, the first coordinate of any resulting
convex combination can be regrouped into the original convex combination of
slices around $x$, and the second coordinate can be regrouped into the one
around $y$.

\section*{Theorem}

\begin{theorem*}
Let $X$ and $Y$ be Banach spaces. If $x\in \operatorname{semiSCS}(B_X)$ and
$y\in \operatorname{semiSCS}(B_Y)$, then
$(x,y)\in \operatorname{semiSCS}(B_{X\oplus_\infty Y})$.
\end{theorem*}

\begin{proof}
Let $\varepsilon>0$. Since $x$ is a semi SCS point of $B_X$, there are slices
$S_i=S(B_X,f_i,\alpha_i)$, $1\leq i\leq n$, and positive weights
$\lambda_i$ with $\sum_i\lambda_i=1$ such that
\[
  \sum_{i=1}^n \lambda_i S_i \subset B(x,\varepsilon).
\]
Similarly, there are slices $T_j=S(B_Y,g_j,\beta_j)$, $1\leq j\leq m$, and
positive weights $\mu_j$ with $\sum_j\mu_j=1$ such that
\[
  \sum_{j=1}^m \mu_j T_j \subset B(y,\varepsilon).
\]
We use the standard convention that slices are defined by nonzero functionals,
and rescale the functionals so that $\|f_i\|=\|g_j\|=1$ for all $i,j$.

Put $Z=X\oplus_\infty Y$. For every pair $(i,j)$ choose
\[
  0<\eta_{ij}<\frac12\min\{\alpha_i,\beta_j\},
\]
and define $h_{ij}\in Z^*=X^*\oplus_1Y^*$ by
\[
  h_{ij}(u,v)=\frac12 f_i(u)+\frac12 g_j(v).
\]
Then $\|h_{ij}\|=1$. Let
\[
  W_{ij}=S(B_Z,h_{ij},\eta_{ij})
\]
be the corresponding slice of $B_Z$. If $(u,v)\in W_{ij}$, then
\[
  \frac12 f_i(u)+\frac12 g_j(v)>1-\eta_{ij},
\]
so
\[
  f_i(u)+g_j(v)>2-2\eta_{ij}.
\]
Because $u\in B_X$, $v\in B_Y$, and $\|f_i\|=\|g_j\|=1$, each of
$f_i(u)$ and $g_j(v)$ is at most $1$. Hence
\[
  f_i(u)>1-2\eta_{ij}>1-\alpha_i,\qquad
  g_j(v)>1-2\eta_{ij}>1-\beta_j.
\]
Thus
\[
  W_{ij}\subset S_i\times T_j.
\]

Consider the convex combination of slices of $B_Z$
\[
  W=\sum_{i=1}^n\sum_{j=1}^m \lambda_i\mu_j W_{ij}.
\]
Let $(u,v)\in W$. Then for some $(u_{ij},v_{ij})\in W_{ij}$,
\[
  (u,v)=\sum_{i=1}^n\sum_{j=1}^m \lambda_i\mu_j (u_{ij},v_{ij}).
\]
For fixed $i$, each $u_{ij}$ belongs to the slice $S_i$, and $S_i$ is convex.
Therefore
\[
  \sum_{j=1}^m \mu_j u_{ij}\in S_i.
\]
Consequently,
\[
  u=\sum_{i=1}^n \lambda_i\left(\sum_{j=1}^m \mu_j u_{ij}\right)
    \in \sum_{i=1}^n \lambda_iS_i\subset B(x,\varepsilon).
\]
The same argument with the roles of $i$ and $j$ reversed gives
\[
  v=\sum_{j=1}^m \mu_j\left(\sum_{i=1}^n \lambda_i v_{ij}\right)
    \in \sum_{j=1}^m \mu_jT_j\subset B(y,\varepsilon).
\]
Hence
\[
  \|(u,v)-(x,y)\|_\infty
  =\max\{\|u-x\|,\|v-y\|\}\leq \varepsilon.
\]
We have found, for arbitrary $\varepsilon>0$, a convex combination of slices
of $B_{X\oplus_\infty Y}$ contained in $B((x,y),\varepsilon)$. Therefore
$(x,y)$ is a semi SCS point of $B_{X\oplus_\infty Y}$.
\end{proof}

\section*{Verification and Novelty Check}

The proof is purely formal from the definitions. The key checks are:
\begin{enumerate}
\item $(X\oplus_\infty Y)^*=X^*\oplus_1Y^*$, so
$h_{ij}=(f_i,g_j)/2$ has norm one when $\|f_i\|=\|g_j\|=1$.
\item A sufficiently thin slice by $h_{ij}$ is contained in the product
$S_i\times T_j$.
\item Slices of a convex set are convex, which justifies the coordinate
regrouping.
\end{enumerate}

The run indexes were searched for arXiv:2602.12699, ``semi SCS'', and
``semi denting''. The only relevant existing row concerns a different
counterexample for arXiv:2307.03621 about semi $w^*$-denting points and
extreme points. Web searches on June 28, 2026, for exact phrases from Question
4.6 and for the source title found only the source arXiv page and no later
answer.

\section*{Scope}

This answers only Question 4.6, the binary $\ell_\infty$-sum converse for
semi SCS points. It does not address the source paper's later questions about
the case $1<p<\infty$ or about converses for projective tensor products.

\section*{References}

\begin{enumerate}
\item S. Basu, P. P. Behera, and S. Seal,
\emph{Semi denting points and related notions in Banach spaces},
arXiv:2602.12699, 2026.
\end{enumerate}

\end{document}
